\chapter{Première mise en service}
\minitoc

\section{Créer le compte super-administrateur}

À la toute première ouverture, la plateforme n'a aucun compte. Elle présente
l'écran \textbf{Première mise en service}, qui ne s'affichera plus ensuite.

\begin{enumerate}[leftmargin=*]
  \item Ouvrir \texttt{http://127.0.0.1:8000}.
  \item Renseigner nom, prénom, nom d'utilisateur, adresse électronique et mot
        de passe.
  \item Valider par \textbf{Créer le compte}.
\end{enumerate}

Le compte obtenu est \emph{super-administrateur} : il a tous les droits, y
compris celui de créer les autres comptes et de transférer son propre rôle.

\begin{tcolorbox}[colback=orange!5!white, colframe=orange!70!black,
  title=\textbf{Attention}, fonttitle=\bfseries, boxrule=0.6pt]
Cet écran n'apparaît qu'une fois, et il n'est protégé par rien. Si la
plateforme est joignable depuis le réseau au moment de sa première mise en
service, la première personne à l'ouvrir en devient l'administrateur. Créez ce
compte immédiatement après le premier démarrage, et de préférence depuis le
poste local.
\end{tcolorbox}

\section{Créer les comptes d'exploitation}

Depuis l'onglet \textbf{Comptes}, accessible à l'administrateur :

\begin{enumerate}[leftmargin=*]
  \item \textbf{Créer un compte}, puis choisir le rôle : analyste, décideur ou
        administrateur.
  \item Un analyste peut aussi s'inscrire lui-même depuis l'écran de connexion~;
        son compte reste alors \emph{en attente} jusqu'à ce qu'un
        administrateur l'admette.
\end{enumerate}

Les mots de passe sont stockés sous forme d'empreintes PBKDF2-HMAC-SHA256 à
600\,000 itérations. Ils ne sont jamais conservés en clair, et ne peuvent pas
être retrouvés — seulement réinitialisés.

\section{Reconnaître l'état du système}

Le bandeau supérieur porte en permanence trois informations. Prenez l'habitude
de les lire avant toute chose.

\begin{description}[leftmargin=3.8cm, style=nextline]
  \item[Autonomie active] La plateforme exécute les gestes qu'elle décide. Si
        la mention est \textbf{Autonomie suspendue}, elle raisonne mais n'agit
        plus : soit un administrateur l'a désactivée, soit le coupe-circuit
        s'est déclenché.
  \item[Mode d'actionnement] \emph{Simulation} ou \emph{production}. En
        simulation, tout est calculé et tracé, mais rien ne touche vos
        équipements. C'est le mode par défaut, et c'est celui des
        démonstrations.
  \item[Compteurs] Le nombre d'interventions ouvertes et de décisions en
        attente, rappelés à côté des onglets concernés.
\end{description}

\section{Premier tour : produire un incident sans risque}

Le meilleur moyen de prendre en main la plateforme est de lui donner une
attaque à traiter. L'onglet \textbf{Démonstration} déclenche des scénarios
réalistes en mode simulation.

\begin{enumerate}[leftmargin=*]
  \item Ouvrir \textbf{Démonstration}.
  \item Choisir un type d'attaque du catalogue --- par exemple \textbf{A5,
        exfiltration de données} --- et cliquer sur son bouton de déclenchement.
  \item Ouvrir \textbf{Vue d'ensemble} : l'intervention y figure, avec les
        gestes que la plateforme vient d'appliquer.
  \item Ouvrir \textbf{Journal d'audit} : chaque étape est inscrite, de la
        réception de l'observation à la notification de l'analyste.
\end{enumerate}

\begin{tcolorbox}[colback=green!4!white, colframe=darkgreen,
  title=\textbf{Bon à savoir}, fonttitle=\bfseries, boxrule=0.6pt]
Les scénarios de démonstration ne fabriquent pas d'attaque : ils fabriquent
\textbf{la charge utile qu'un collecteur émettrait} pour l'attaque décrite. La
plateforme ne sait pas qu'elle est en démonstration --- elle traite ces
observations exactement comme celles d'une sonde réelle.
\end{tcolorbox}

\section{Remise à zéro}

L'onglet \textbf{Démonstration} propose une remise à zéro qui efface les
interventions et les actions produites par les scénarios. Elle \textbf{ne
touche pas} au journal d'audit : celui-ci est immuable par construction, et
l'effacer reviendrait à détruire la seule preuve de ce que le système a fait.
